\section{Discussion and Limitations}
\label{sec:discussion}

\subsection{What the Analytical Model Does---and Does Not---Show}

Equation~\eqref{eq:erlang-recursion} is useful because it converts capacity and a desired loss probability into a transparent offered-load budget. It does not prove that a finite trace, an MMPP, correlated retries, or a queueing LLM subscriber will meet that probability. Similarly, numerically selecting $A$ so that $B_n(A)=\alpha$ and then plotting $B_n(A)$ is an identity, not validation. Our evidence instead comes from correlated completion, expiry, task accuracy, PUBACK timing, and blocking observed outside that equation.

The experiments also show why a single objective is hazardous. The analytical and adaptive ablations sometimes approach or beat the blocking target by setting a deadline that removes admitted state before useful work completes. A deployment must decide whether rejecting at admission, expiring after admission, or returning a lower-quality fast result has the lowest cost. \method{} currently emphasizes completion of admitted work through feasibility floors; this can raise blocking. A future formulation should optimize a weighted utility or constrained vector of blocking, expiry, quality, and class fairness rather than treating one target as sufficient.

The conventional comparators reinforce this point. An occupancy threshold is difficult to beat on simplicity and, under MMPP traffic, converts burst pressure into early rejection with almost no post-admission expiry. Static reservation can protect admission opportunity for the minority CPU-preferred class, but the protected jobs may still expire and their occupancy can displace many shorter GPU-preferred jobs. The combined threshold/reservation rule is therefore not automatically safer: in this workload it is conservative for the majority class without adapting to the protected class's service drift. \method{} is most useful when service and load move together, while a fixed threshold remains a credible choice when early rejection is acceptable and the workload is persistently bursty.

\subsection{Deployment Scope}

The controller is an application-layer sidecar, not a modified Mosquitto binary. This choice improves portability and preserves MQTT conformance, but adds a process and duplicates selected request metadata. A production design could move the same state machine into a broker plugin, persistent gateway service, or replicated control plane. Crash recovery, exactly-once application effects, multi-broker correlation, authentication/authorization, TLS overhead, and adversarial correlation identifiers are outside this prototype. MQTT QoS~1 can redeliver an Application Message; a production controller must make $id_i$ admission and completion idempotent.

The experiment uses loopback networking and one physical host. It isolates broker semantics, execution-class differences, and LLM service time, but does not reproduce wireless loss, WAN delay, battery/thermal throttling, or many-device contention. The CPU-only endpoint is a software-enforced execution class on the Ryzen host. The reported process memory is therefore an implementation measurement, not a general broker-overhead model.

The 240-example corpus provides exact, reproducible labels and makes conflicting context controllable, but it remains synthetic. It spans two prompt templates, four sensor domains, and direct-reading, history-dependent, and unstructured-log requests; half of the items are held out. Profiles across three Qwen2.5 sizes and Llama~3.1--8B expose model dependence, but broader domain datasets, quantizations, and human-authored logs remain necessary. Strict JSON validity is relevant to machine-to-machine IoT control, yet other applications require F1, calibration, factuality, human preference, or safety review. Quality profiles must be regenerated after any model, runtime, or policy change.

The p95 FIFO floor in~\eqref{eq:floors}, feedback gain, epoch duration, and 30\% rate limiter were engineering choices evaluated on one host. They are disclosed to enable replication but are not claimed optimal. A quantile predictor with conformal calibration or online uncertainty could replace the empirical bound. Likewise, the linear queue prediction assumes one serial worker per class and ignores batching. Datacenter engines such as vLLM, Sarathi-Serve, DistServe, and NanoFlow have different service dynamics~\cite{kwon2023vllm,agrawal2024sarathi,zhong2024distserve,zhu2025nanoflow}.

\subsection{Arrival Diagnostics and Robust Control}

One-second dispersion separates the Poisson and MMPP runs used here, but it can miss sub-second synchronized batches: a process may place several events at nearly the same instant while keeping one-second counts regular. A deployment should compute dispersion and autocorrelation at multiple bin widths, track interarrival coefficient of variation, and detect rate changes. The controller logs one-second dispersion and each MMPP state transition; every 120-s live MMPP schedule is checked for visits to both states.

The MMPP results delimit the operating region of mean-rate control. Neither empirical calibration nor feedback makes a mean-rate Erlang budget robust to bursts that fill the store within an epoch. The replay therefore includes a token-bucket front end with a bounded two-second shaping delay and reports broker blocking, shaper drops, total admission loss, and shaping delay separately. This distinction prevents an apparent blocking reduction from hiding loss at the shaper. Predictive change detection, risk-sensitive control, and synchronized-retry suppression remain open. If records consume different numbers of memory units, a multirate Kaufman--Roberts calculation or direct measured admission control is appropriate.

\subsection{Artifact and Statistical Scope}

The artifact at \artifacturl{} records raw JSONL inference profiles, per-run configurations and events, control-epoch CSV files, broker logs, simulation seeds, aggregate tables, and figure-generation code. The 310 selected live runs use ten paired seeds per main comparison so policies see identical arrival schedules. Student-$t$ intervals summarize between-run variability, while two-sided exact sign-flip tests operate on paired blocking, expiry, and goodput differences. Conventional policies are order-balanced within their campaign, but the \method{} and conventional campaigns occurred in separate time blocks; pairing removes input-trace variation, not every possible long-term host-state effect. Replay reports 95\% intervals over at least 100 seeds for each main point. Because replay resamples a finite empirical profile, it quantifies variability conditional on this measured host/model distribution; it does not create evidence for unmeasured devices.

\section{Conclusion}
\label{sec:conclusion}

This paper presents an implemented, standards-consistent application-state mechanism for MQTT--LLM pipelines. MQTT delivery acknowledgments proceed normally; LLM completion is a separate correlated MQTT~5 message, and a co-located controller holds bounded state until completion or expiry. The deadline controller uses Erlang--B as a conditional feedforward budget, converts that budget with an empirical truncated-mean inversion, and adds measured feedback, class feasibility floors, and load-aware routing. The subscriber selects a profiled inference configuration by task utility under a tail-latency constraint, with no requirement to fill the context window.

The Mosquitto/Qwen2.5 study shows that millisecond protocol ACK latency coexists with second-scale application completion. Class-conditioned selection improves live accuracy by 12.33 percentage points over compact selection but remains below maximum-context and token-max baselines. \method{} has lower blocking than occupancy-threshold admission, static class reservation, and their combination in all three traffic conditions. During simultaneous arrival-rate and model-service drift it also gains 0.527/s over the strongest conventional comparator and 1.273/s over the combined baseline. Under MMPP input, however, occupancy-threshold admission gains 0.186/s by blocking more requests early and nearly eliminating expiry, and the mean-rate controller can also underperform fixed-deadline control. Protocol-correct completion, measured quality/latency profiles, exact paired inference, and explicit feasibility checks make both the changing-load advantage and the burst limitation observable.

\section*{Acknowledgment}

This research was supported by the ANCHOR program through the Ulsan ANCHOR Center, funded by the Ministry of Education (MOE) and the Ulsan Metropolitan City, Republic of Korea (2026-ANCHOR-07-001).
